Because the frequency reports on culture talk are only available in the form of vague quantifiers we need a way to ``norm'' the individual's response to these items to hold constant idiosyncratic interpretations of the frequency categories.  The aim is to determine whether persons in certain socio-demographic groups are more likely to report higher frequencies of culture talk than others, but the problem is that no ``objective'' estimate of frequency is available. My strategy is to use the self-reported frequencies on culture consumption as a baseline with which to determine whether persons engage in more culture talk than would be expected \emph{in relation to their own self-reported consumption frequencies in that cultural domain.} I proceed by arranging the data in the form of matched-pairs \citep[247-256]{agresti07}.  Each individual's response to the culture-talk item is matched against their own response to the consumption item (e.g. serving as their own ``control''), contributing two observations to the analysis.  An example of this arrangement is shown in table~\ref{tab:matched} for three hypothetical cases.   Accordingly, we can think of the (joint) outcome of interest as the subject-specific cross-classification of the response to the frequency of consumption item versus their response to the frequency of cultural talk item in that cultural domain, both expressed in terms of vague-quantifiers.  The data (collapsed across subjects in the form of a square four by four table) are shown in table~\ref{tab:contingency}.  The matched pairs design also has the advantage of providing a simple and parsimonuous way of allowing me to test the hypothesis of whether the effect of different socio-demographic predictors is the same or different across the two outcomes \citep[Pp.]{agresti07}.

Nevertheless, even after matching the each person's subjective estimate of their consumption and talk frequency to generate a subject-specific deviation from what they perceive to be their typical behavior as an estimate of the frequency of culture-talk, it is still possible that there exists structured heterogeneity in the way that such terms as ``very often'' are interpreted \emph{across} respondents in different social locations. I deal with the problem of cross-individual variation in the interpretation of the cutpoints as follows.  Rather than estimating a standard ordered logistic regression model, which restricts the cutpoint estimates to be the same for all respondents, I specify a \emph{generalized ordered logit model} \citep{williams06}. This model relaxes the assumption of cross-individual homogeneity in the cutpoint estimates to allow for cross-individual heterogeneity--based on the available socio-demographic covariates---in the location of the cutpoints.  With this model, I am able to estimate the joint socio-demographic variation in the propensity to engage in culture-talk in relation to culture-consumption and in the propensity to match different objective frequencies to the frequency labels provided in the survey item. The analyses shown in the Appendix, which compare objective to subjective report frequencies for the consumption items show that there is predictable socio-demographic heterogeneity in the interpretation of what constitutes ``very often'' versus ``fairly often.''  In essence, those socio-demographic groups---the educated, young and wealthy---who are more likely to consume culture more frequently have ``tougher'' thresholds (shifted to the right of the distribution) as to what constitutes ``very often.''  Relatedly, those socio-demographic groups who are less likely to be culture consumers (older, less educated respondents) tend to think of lower objective frequencies as ``very often'' in comparison to their younger and more privileged counterparts.  We should expect similar effects to obtain in the analysis reported below.

The model to be estimated is one where $logit[P(Y_{i1}\leq j)]=\alpha_{ij}$, and $logit[P(Y_{i2}\leq j)]=\alpha_{ij} + \beta$ \citep[254]{agresti07}, where $\beta$ is the subject-specific difference in the probability of engaging in talk over consumption:

\begin{equation}
logit[P(Y_{it}\leq j)]=\alpha_{ij} + \beta_t, t=1,2
\label{eq:agresti1}
\end{equation}•

Here I use the ordered structure of the responses by restricting the comparisons to logits of cumulative probabilities (e.g. estimating a fixed-effects ordered logit model).  The dependent variable is the---identically coded---joint response to the culture consumption and culture talk items ($Y_{it}$) for individual.  Here $t=1$ for the consumption item and $t=2$ for the cultural talk item respectively.  The indicator variable $\beta_t$ can be interpreted as giving us the difference in response probabilities between the talk and consumption items for each person.   If the parameter is positive ($\beta_t>0$) this indicates that culture-talk is more frequent in that cultural domain in comparison to reported frequencies of consumption;  $\beta_t<0$ indicates the reverse: consumption is more prevalent than talk in that domain.  Finally if $\beta_t \approx 0$, then than means that both consumption and talk are equally prevalent. Note that the $\alpha_{ij}$ condition on all idiosyncratic variation due to cross-individual differences ($i$) in the interpretation of the cutoff points ($j$).

If we are interested in cross-group differences in the likelihood of engaging in culture talk (relative to culture consumption) we can introduce interaction effects between socio-demographic covariates and the indicator variable $\beta$.  Here the individual level covariate \emph{modify} the baseline propensity of some genres to feature an imbalance of talk over consumption (or vice versa).  Thus, for individual-level covariates, $X_1$, $X_2$, $X_3${\ldots}$X_k$, Equation~\ref{eq:agresti1}, generalizes to:

\begin{equation}
logit[P(Y_{it}\leq j)]=\alpha_{ij} + \beta_t+\sum_k\delta_kX_{ikt}, t=1,2
\label{eq:agresti1}
\end{equation}•

Where $X_{ikt}=\beta_t \times X_{ik}$ for all $k$ variables.  